\chapter{Hardware Architecture and Implementation}

\section{Introduction}

This chapter describes the physical hardware platform used to implement the firmware architecture defined in previous chapters. 

The objective is to:

\begin{itemize}
	\item Describe the electronic components used
	\item Justify hardware selection decisions
	\item Explain system-level interconnections
	\item Demonstrate consistency with the hardware abstraction strategy
\end{itemize}

The hardware platform was designed to support modularity, expandability, and reliable real-time operation.

% ------------------------------------------------

\section{System Overview}

The hardware system consists of:

\begin{itemize}
	\item A primary microcontroller unit (MCU)
	\item Wireless communication interface
	\item Motor driver stage
	\item Servo control outputs
	\item Sensor interfaces
	\item GPIO expansion module
	\item Power management circuitry
\end{itemize}

All subsystems are electrically integrated while maintaining logical separation consistent with the firmware structure.

% ------------------------------------------------

\section{Microcontroller Unit}

The core processing element is the \textbf{ESP32}

The ESP32 was selected due to:

\begin{itemize}
	\item Integrated Bluetooth capability
	\item Dual-core processing architecture
	\item Hardware PWM peripherals
	\item I2C and UART support
	\item Sufficient RAM and Flash resources
\end{itemize}

Its integrated wireless stack eliminates the need for external communication modules, reducing system complexity.

% ------------------------------------------------

\section{Motor Driver Stage}

DC motor actuation is achieved through an H-bridge driver module.

The driver stage provides:

\begin{itemize}
	\item Bidirectional current flow
	\item PWM speed control
	\item Electrical isolation between logic and motor supply
\end{itemize}

The driver was selected based on:

\begin{itemize}
	\item Current handling capacity
	\item Voltage compatibility
	\item Switching efficiency
\end{itemize}

The logic-level control inputs are directly interfaced with MCU PWM outputs.

% ------------------------------------------------

\section{Servo Control Interface}

Steering and auxiliary actuation use standard hobby servo motors.

Servo signals are generated via:

\begin{itemize}
	\item Dedicated hardware PWM channels
	\item 50 Hz refresh frequency
	\item Pulse width range typically between 1 ms and 2 ms
\end{itemize}

Because servo timing is hardware-assisted, software timing jitter does not affect signal stability.

% ------------------------------------------------

\section{Sensor Subsystem}

Sensors integrated into the system include:

\begin{itemize}
	\item Analog voltage measurement inputs
	\item Digital state inputs
	\item I2C-based sensor devices
\end{itemize}

Analog signals are read through the MCU ADC channels, while digital expansion is handled via an external GPIO expander.

This hybrid approach increases scalability without exhausting native MCU pins.

% ------------------------------------------------

\section{GPIO Expansion Module}

An external I2C-based GPIO expander is used to extend digital input/output capacity.

Advantages:

\begin{itemize}
	\item Reduced pin usage on MCU
	\item Modular hardware scaling
	\item Electrical separation of subsystems
\end{itemize}

The I2C bus operates at a standard clock frequency compatible with reliable real-time operation.

% ------------------------------------------------

\section{Power Distribution Architecture}

The system includes separate supply domains:

\begin{itemize}
	\item Logic supply (3.3V)
	\item Servo supply (typically 5V)
	\item Motor supply (dependent on motor rating)
\end{itemize}

Separation of power domains prevents:

\begin{itemize}
	\item Voltage dips during motor startup
	\item Electrical noise propagation
	\item Microcontroller instability
\end{itemize}

Decoupling capacitors are placed near high-current devices to reduce transient disturbances.

% ------------------------------------------------

\section{Signal Integrity Considerations}

The following measures were taken to preserve signal integrity:

\begin{itemize}
	\item Short PWM trace lengths
	\item Common ground reference
	\item Shielded or twisted wiring for motor leads
	\item Pull-up resistors for I2C lines
\end{itemize}

These precautions reduce electromagnetic interference (EMI) and improve system reliability.

% ------------------------------------------------

\section{Fail-Safe and Protection Mechanisms}

Hardware-level protection includes:

\begin{itemize}
	\item Reverse polarity protection (if implemented)
	\item Current limitations via driver module
	\item Brown-out detection (MCU internal feature)
\end{itemize}

These mechanisms enhance operational safety and prevent hardware damage.

% ------------------------------------------------

\section{Hardware–Firmware Consistency}

The hardware implementation directly reflects the abstraction strategy defined in Chapter 2:

\begin{itemize}
	\item All pins are centrally defined
	\item Peripheral interfaces match modular firmware drivers
	\item No hard-coded electrical assumptions exist in application logic
\end{itemize}

This ensures that hardware changes (e.g., driver replacement) require minimal firmware modification.

% ------------------------------------------------

\section{Scalability and Expansion Potential}

The hardware platform supports future extensions such as:

\begin{itemize}
	\item Additional sensors
	\item Higher current motor stages
	\item External communication interfaces
	\item Advanced feedback control modules
\end{itemize}

Expansion is supported through both available MCU peripherals and modular firmware structure.

% ------------------------------------------------

\section{Summary}

The hardware architecture provides:

\begin{itemize}
	\item Integrated wireless capability
	\item Deterministic actuator control
	\item Modular expansion support
	\item Electrical stability and protection
	\item Direct alignment with firmware abstraction principles
\end{itemize}

This cohesive hardware–software integration demonstrates a complete embedded system implementation rather than isolated firmware development.